%% ====================================================================
\section{Theoretical Framework}
\label{sec:theory}
%% ====================================================================

\subsection{Dynamical Quantum Phase Transitions}
\label{sec:dqpt}

Dynamical quantum phase transitions (DQPTs) were introduced by Heyl, Polkovnikov, and
Kehrein~\cite{heyl2013} as real-time analogs of equilibrium phase transitions. The
central object is the Loschmidt amplitude:

\begin{definition}[Loschmidt amplitude]\label{def:loschmidt}
Given an initial state $|\psi_0\rangle$ (typically the ground state of a pre-quench
Hamiltonian~$H_0$) and a post-quench Hamiltonian~$H$, the Loschmidt amplitude is
\begin{equation}\label{eq:loschmidt}
  \mathcal{G}(t) = \langle\psi_0| e^{-iHt} |\psi_0\rangle.
\end{equation}
The return rate (or rate function) is defined as
\begin{equation}\label{eq:rate-general}
  g(t) = -\frac{1}{N}\ln|\mathcal{G}(t)|^2,
\end{equation}
where $N$ is the system size. A DQPT occurs at time~$t^*$ if $g(t)$ is non-analytic
at~$t = t^*$.
\end{definition}

The connection to equilibrium statistical mechanics arises through the formal
correspondence $t \to -i\beta$ (Wick rotation), under which the Loschmidt amplitude
becomes
\begin{equation}\label{eq:wick}
  \mathcal{G}(-i\beta) = \langle\psi_0| e^{-\beta H} |\psi_0\rangle.
\end{equation}
For the specific case where $|\psi_0\rangle$ is a uniform superposition over energy
eigenstates, this equals $Z(\beta)/\dim(\mathcal{H})$. In general, the relation depends
on the overlap of $|\psi_0\rangle$ with the eigenstates of~$H$.

Non-analyticities in the free energy density $f(\beta) = -(1/N)\ln Z(\beta)$ as a
function of inverse temperature correspond to equilibrium phase transitions. By the
Lee--Yang theorem~\cite{leeyang1952}, these non-analyticities are controlled by the
zeros of the partition function in the complex $\beta$-plane. Analogously, DQPTs are
controlled by the zeros of the Loschmidt amplitude in the complex time plane. When
these zeros cross the real time axis, the rate function $g(t)$ exhibits non-analytic
behavior (cusps, divergences, or discontinuities in derivatives), signaling a
DQPT~\cite{heyl2013}.

\paragraph{Connection to Lee--Yang zeros.}
In the Lee--Yang theory~\cite{leeyang1952,fisher1965}, equilibrium phase transitions
arise when zeros of the partition function $Z(\beta,h)$---viewed as a function of
complex fugacity or magnetic field---pinch the real axis in the thermodynamic limit.
The density of these zeros determines the order and nature of the transition. In the
DQPT framework, the analogous role is played by the Fisher zeros of the Loschmidt
amplitude in the complex time plane. The Wei et~al.\ construction arranges for these
Fisher zeros to coincide with the nontrivial zeros of the Riemann zeta function.


\subsection{The Wei et al.\ Construction}
\label{sec:wei}

\subsubsection{System~1: Probe Spin Coherence with Logarithmic Spectrum}
\label{sec:system1}

\paragraph{Setup.}
Consider a composite system consisting of a probe spin-$\tfrac{1}{2}$ coupled to a
bath of $N$ quantum levels. The bath Hamiltonian has a logarithmic energy spectrum:
\begin{equation}\label{eq:H0}
  H_0 = \sum_{n=1}^{N} \ln(n)\,|n\rangle\langle n|,
\end{equation}
where $\{|n\rangle\}_{n=1}^{N}$ is an orthonormal basis for the bath Hilbert space.

\begin{remark}[On the spectrum]\label{rem:spectrum}
The logarithmic spectrum $E_n = \ln(n)$ is reverse-engineered to produce the Dirichlet
series in the partition function; it is not a spectrum that arises naturally in a
standard condensed-matter or atomic physics context. However, logarithmic spectra do
arise in the theory of quantum mechanics on hyperbolic surfaces: the Selberg zeta
function
\begin{equation}\label{eq:selberg}
  Z_\Gamma(s) = \prod_{\{p\}} \prod_{k=0}^{\infty}
  \bigl(1 - e^{-(s+k)\ell_p}\bigr),
\end{equation}
associated to a discrete subgroup $\Gamma$ of $\mathrm{SL}(2,\R)$, encodes lengths
$\ell_p$ of primitive closed geodesics, and the corresponding quantum Hamiltonian (the
Laplace--Beltrami operator on the quotient surface) has a spectral theory intimately
related to zeta functions~\cite{edwards1974}. This provides a potential natural
realization of the logarithmic spectrum, though the precise connection to the
Wei et~al.\ construction remains to be established.
\end{remark}

The probe spin couples to the bath such that the effective Hamiltonian for the coupled
evolution is
\begin{equation}\label{eq:H-coupled}
  H = \sigma_z \otimes H_0,
\end{equation}
where $\sigma_z$ is the Pauli-$Z$ operator acting on the probe spin.

\paragraph{Thermal state preparation.}
The bath is prepared in a thermal (Gibbs) state at inverse temperature~$\beta$:
\begin{equation}\label{eq:thermal}
  \rho_{\mathrm{bath}}(\beta) = \frac{e^{-\beta H_0}}{Z(\beta,H_0)},
\end{equation}
where the partition function is
\begin{equation}\label{eq:partition}
  Z(\beta,H_0) = \Tr\bigl[e^{-\beta H_0}\bigr]
  = \sum_{n=1}^{N} e^{-\beta\ln(n)}
  = \sum_{n=1}^{N} n^{-\beta}.
\end{equation}
Note that $Z(\beta,H_0)$ is precisely the partial zeta function. In the thermodynamic
limit $N \to \infty$ and for $\beta > 1$, this converges to
$\zeta(\beta)$~\cite{titchmarsh1986}.

The probe spin is prepared in the state
$|{+}\rangle = (|{\uparrow}\rangle + |{\downarrow}\rangle)/\sqrt{2}$.

\paragraph{Time evolution.}
The probe-bath system is prepared in the product state
\begin{equation}\label{eq:initial-state}
  \rho_{\mathrm{total}} = |{+}\rangle\langle{+}| \otimes \rho_{\mathrm{bath}}(\beta),
\end{equation}
where $|{+}\rangle\langle{+}|$ is the pure state projector for the probe spin and
$\rho_{\mathrm{bath}}(\beta)$ is the thermal state defined above. Under the coupled
Hamiltonian $H = \sigma_z \otimes H_0$, the total state evolves as
\begin{equation}\label{eq:evolution}
  \rho_{\mathrm{total}}(t) = e^{-iHt}\,\rho_{\mathrm{total}}\,e^{iHt}.
\end{equation}
The reduced density matrix of the probe spin, obtained by tracing over the bath
(partial trace), yields the coherence (off-diagonal element):
\begin{equation}\label{eq:coherence}
  \langle\sigma_+\rangle(t) = \Tr_{\mathrm{bath}}\bigl[\rho_{\mathrm{bath}}\,e^{2iH_0 t}\bigr].
\end{equation}

\paragraph{Derivation of the accumulated phase factor.}
We compute:
\begin{align}\label{eq:coherence-deriv}
  \langle\sigma_+\rangle(t)
  &= \Tr\bigl[\rho_{\mathrm{bath}}\,e^{2iH_0 t}\bigr] \nonumber\\
  &= \frac{1}{Z(\beta)} \sum_{n=1}^{N} e^{-\beta\ln(n)}\,e^{2i\ln(n)\,t} \nonumber\\
  &= \frac{1}{Z(\beta)} \sum_{n=1}^{N} n^{-\beta}\,n^{2it} \nonumber\\
  &= \frac{1}{Z(\beta)} \sum_{n=1}^{N} n^{-(\beta - 2it)} \nonumber\\
  &= \frac{S_N(\beta - 2it)}{Z(\beta)},
\end{align}
where $S_N(s) = \sum_{n=1}^{N} n^{-s}$ is the partial Dirichlet series.

Note that the coherence $\langle\sigma_+\rangle$ involves $S_N(\beta - 2it)$, not
$S_N(\beta + it)$. Wei et~al.~\cite{wei2025} introduce three modifications to obtain
the accumulated phase factor $\mathcal{L}(\beta,t)$:
\begin{enumerate}
  \item[(a)] A reparametrization $t \to t/2$ absorbs the factor of~2: replacing $2t$
    by~$t$ in the exponent gives $S_N(\beta - it)$. This is purely a rescaling of the
    time variable and does not affect the location of zeros (a zero of $S_N(\beta - it_0)$
    at time~$t_0$ corresponds to a zero of $S_N(\beta - 2it_0')$ at time $t_0' = t_0/2$).
  \item[(b)] A sign convention: the coherence involves $n^{-(\beta - it)} =
    n^{-\beta}\,n^{it}$, while the standard Dirichlet series uses
    $n^{-(\beta + it)} = n^{-\beta}\,n^{-it}$. The passage from $-it$ to $+it$
    corresponds to complex conjugation of the time evolution phase, equivalent to
    choosing the opposite sign convention for the coupling Hamiltonian. This does not
    affect the zero locations (if $\eta(\beta + it_0) = 0$, then $\eta(\beta - it_0) = 0$
    by the reflection principle, since the eta coefficients are real).
  \item[(c)] An alternating sign $(-1)^{n+1}$ is introduced, corresponding physically to
    a staggered coupling between the probe and even/odd bath levels. This converts the
    ordinary Dirichlet series into the alternating Dirichlet eta series.\footnote{This
    staggered coupling requires that the probe spin distinguishes the parity of the
    energy index~$n$ and reverses its coupling accordingly. Since $n$ is an energy index
    rather than a spatial coordinate, this coupling is \emph{non-local in energy space}
    and does not arise naturally in any known many-body system. Together with the
    logarithmic spectrum $E_n = \ln n$, this reinforces the point made in
    Section~\ref{sec:scope}: the quantum system is \emph{engineered} to reproduce the
    Dirichlet eta function, and the ``physics'' is a mathematical reformulation rather
    than an independent physical mechanism.}
\end{enumerate}
With these three modifications, the accumulated phase factor becomes:
\begin{equation}\label{eq:L-def}
  \mathcal{L}(\beta,t) = \frac{1}{Z(\beta)}
  \sum_{n=1}^{N} (-1)^{n+1}\,n^{-(\beta+it)}.
\end{equation}
The factor $(-1)^{n+1}$ converts the ordinary Dirichlet series into the alternating
Dirichlet series, which defines the Dirichlet eta function:
\begin{equation}\label{eq:eta-def}
  \eta(s) = \sum_{n=1}^{\infty} (-1)^{n+1}\,n^{-s}
  = (1 - 2^{1-s})\,\zeta(s).
\end{equation}

\paragraph{Why $\mathcal{L}$ vanishes at zeta zeros.}
We now prove the key result. In the thermodynamic limit $N \to \infty$:

\medskip\noindent
\textbf{Step~1.} The numerator of $\mathcal{L}(\beta,t)$ converges:
\begin{equation}\label{eq:eta-convergence}
  \sum_{n=1}^{\infty} (-1)^{n+1}\,n^{-(\beta+it)} = \eta(\beta+it).
\end{equation}
This series converges for $\re(\beta+it) = \beta > 0$ by Dirichlet's test (Abel
summation): the partial sums of $(-1)^{n+1}$ are bounded and $n^{-\beta}$ decreases
monotonically to zero~\cite{titchmarsh1986}. (Note: the classical Leibniz alternating
series test applies only to real series; for complex $s = \beta + it$ with $t \ne 0$,
convergence follows from Dirichlet's test.)

\medskip\noindent
\textbf{Step~2.} The Dirichlet eta function satisfies the identity
\begin{equation}\label{eq:eta-zeta}
  \eta(s) = (1 - 2^{1-s})\,\zeta(s)
\end{equation}
for all $s \in \C$ (by analytic continuation). This identity is proved by observing
that for $\re(s) > 1$:
\begin{equation}\label{eq:eta-zeta-proof}
  \zeta(s) - \eta(s) = 2\sum_{n=1}^{\infty}(2n)^{-s} = 2^{1-s}\,\zeta(s),
\end{equation}
whence $\eta(s) = \zeta(s) - 2^{1-s}\zeta(s) = (1-2^{1-s})\zeta(s)$. Both sides are
meromorphic on~$\C$ and agree on $\re(s) > 1$; by the identity theorem for meromorphic
functions, they agree everywhere on~$\C$.

\medskip\noindent
\textbf{Step~3.} The prefactor $(1 - 2^{1-s})$ vanishes when $2^{1-s} = 1$, i.e., when
$(1-s)\ln 2 = 2\pi ik$ for some integer~$k$. Writing $s = \beta + it$:
\begin{equation}\label{eq:prefactor-eq}
  (1 - \beta - it)\ln 2 = 2\pi ik.
\end{equation}
Separating real and imaginary parts:
\begin{align}
  \text{Real:}\quad & (1-\beta)\ln 2 = 0, \quad\text{hence } \beta = 1. \label{eq:pf-real}\\
  \text{Imaginary:}\quad & -t\ln 2 = 2\pi k, \quad\text{hence }
    t = -2\pi k/\ln 2 = 2\pi(-k)/\ln 2. \label{eq:pf-imag}
\end{align}
Since $k$ ranges over all integers, the set of solutions is $t = 2\pi m/\ln 2$ for
$m \in \Z$ (relabeling $m = -k$). The prefactor zeros thus lie at
$s = 1 + i(2\pi m/\ln 2)$ for $m \in \Z$. Since $\beta = 1$ lies \emph{outside} the
open critical strip $0 < \beta < 1$ (and no nontrivial zeta zero has $\re(s) = 1$
by the prime number theorem~\cite{titchmarsh1986}), these prefactor zeros never
coincide with nontrivial zeta zeros.

\medskip\noindent
\textbf{Step~4.} Therefore, for $s = \beta + it$ with $0 < \beta < 1$:
\begin{equation}\label{eq:eta-zero-equiv}
  \eta(s) = 0 \;\Longleftrightarrow\; \zeta(s) = 0.
\end{equation}

\medskip\noindent
\textbf{Step~5: The partition function divergence.}
The partition function $Z(\beta) = \sum_{n=1}^{N} n^{-\beta}$ converges to
$\zeta(\beta)$ only for $\beta > 1$. For $\beta = \tfrac{1}{2}$ (the critical case
for~RH), $Z(\tfrac{1}{2}) = \sum 1/\sqrt{n}$ diverges as $2\sqrt{N}$ by the integral
test~\cite{titchmarsh1986}.

This means the naive ratio $\mathcal{L} = \eta / Z$ is a $0/\infty$ form at the zeta
zeros when $\beta = \tfrac{1}{2}$. We must handle this carefully.

Following Wei et~al.~\cite{wei2025}, the physically meaningful quantity is the
\emph{finite-$N$} expression:
\begin{equation}\label{eq:LN-def}
  \mathcal{L}_N(\beta,t) = \frac{\sum_{n=1}^{N}(-1)^{n+1}n^{-(\beta+it)}}
  {\sum_{n=1}^{N} n^{-\beta}}.
\end{equation}
For finite~$N$, both numerator and denominator are finite. The key observation is that
the \emph{rate function} (free energy density) has a well-defined thermodynamic limit:
\begin{equation}\label{eq:F1-def}
  \mathcal{F}_1(\beta,t) = -\lim_{N\to\infty} \frac{1}{\ln N}\,
  \ln|\mathcal{L}_N(\beta,t)|.
\end{equation}
To see this, note that for the numerator:
\begin{itemize}
  \item At a zeta zero $s_0 = \beta_0 + it_0$: the partial eta sum
    $S_N^{\eta}(s_0) = O(N^{-\beta_0})$ by the Dirichlet test remainder
    bound~\cite{titchmarsh1986}.
  \item Away from zeros: $S_N^{\eta}(s) \to \eta(s) \ne 0$, a nonzero constant.
\end{itemize}
For the denominator $Z_N(\beta)$:
\begin{itemize}
  \item For $0 < \beta < 1$: $Z_N(\beta) \sim N^{1-\beta}/(1-\beta)$ by
    Euler--Maclaurin summation~\cite{titchmarsh1986}.
\end{itemize}
Therefore:
\begin{equation}\label{eq:LN-scaling}
  |\mathcal{L}_N(\beta,t)| \sim |S_N^{\eta}(s)|\,\frac{1-\beta}{N^{1-\beta}}.
\end{equation}
At a zero ($s = s_0$): $|S_N^{\eta}(s_0)| = O(N^{-\beta_0})$, so
$|\mathcal{L}_N| \sim N^{\beta-\beta_0-1}$.
Away from a zero: $|S_N^{\eta}(s)| = O(1)$, so $|\mathcal{L}_N| \sim N^{\beta-1}$.

We verify the limit~\eqref{eq:F1-def} exists. Away from zeros,
$\ln|\mathcal{L}_N| = \ln|\eta(s)(1-\beta)| - (1-\beta)\ln N$, giving:
\begin{align}
  \mathcal{F}_1 &= -\lim_{N\to\infty}\frac{1}{\ln N}
    \bigl[\ln|\eta(s)(1-\beta)| - (1-\beta)\ln N\bigr] \nonumber\\
  &= -\lim_{N\to\infty}\frac{\ln|\eta(s)(1-\beta)|}{\ln N} + (1-\beta)
    = 1-\beta. \label{eq:F1-away}
\end{align}
At a zero, $\ln|\mathcal{L}_N| = \ln D + (\beta-\beta_0-1)\ln N$:
\begin{align}
  \mathcal{F}_1 &= -\lim_{N\to\infty}\frac{1}{\ln N}
    \bigl[\ln D + (\beta-\beta_0-1)\ln N\bigr] \nonumber\\
  &= \beta_0 + 1 - \beta. \label{eq:F1-zero}
\end{align}
Both limits exist and are finite. This confirms that $\mathcal{F}_1$ is well-defined
despite the divergence of $Z(\beta)$ for $\beta \le 1$.

\medskip\noindent
\textbf{Step~6: Equivalence statement.} The Riemann Hypothesis is equivalent to:
\begin{equation}\label{eq:rh-dqpt-equiv}
  \text{For all } t \in \R:\;
  \mathcal{F}_1(\beta,t) \text{ exhibits a singularity (DQPT) only at }
  \beta = \tfrac{1}{2}.
\end{equation}
A counterexample to~RH---a zero $\rho = \beta_0 + it_0$ with
$\beta_0 \ne \tfrac{1}{2}$---would manifest as a singularity in $\mathcal{F}_1$ at
$\beta = \beta_0$. This is numerically detectable by scanning $\mathcal{F}_1(\beta,t)$
over a grid in the $(\beta,t)$ plane and looking for anomalous values of
$\mathcal{F}_1$ at $\beta \ne \tfrac{1}{2}$.

\paragraph{Convergence of the eta partial sums.}
The Dirichlet eta series converges for $\re(s) > 0$ by Dirichlet's test (Abel
summation), but only conditionally (not absolutely) for $0 < \re(s) \le 1$. The
convergence rate is:
\begin{equation}\label{eq:eta-remainder}
  |\eta(s) - \eta_N(s)| \le C(s)\,(N+1)^{-\beta},
\end{equation}
where $C(s) = O(1)$ for $s$ bounded away from $\re(s) = 0$, but $C(s)$ grows as
$|\im(s)|$ increases for fixed $\re(s)$ near~0. This is $O(N^{-1/2})$ at
$\beta = \tfrac{1}{2}$, requiring $N \sim 10^{2k}$ to achieve $k$ digits of precision.
For double precision (${\sim}15$ digits), we need $N \sim 10^{30}$, which is infeasible
by direct summation. We address acceleration techniques in Section~\ref{sec:methods}.

\paragraph{Numerical acceleration.}
To circumvent the slow convergence, we use:
\begin{enumerate}
  \item[(a)] The Borwein--Cohen--Villegas--Zagier acceleration~\cite{borwein2000}:
    transforms the conditionally convergent eta series into a geometrically convergent
    series with error $O(3^{-N})$, requiring only $N \sim 40$ terms for double precision.
  \item[(b)] The Euler--Maclaurin summation formula applied to the non-alternating sum
    $S_N(s)$, combined with the identity $\eta(s) = (1-2^{1-s})\zeta(s)$ and the
    Riemann--Siegel formula for~$\zeta$.
\end{enumerate}
Both methods are detailed in Section~\ref{sec:methods}.

\begin{remark}[Trivial zeros]\label{rem:trivial}
The trivial zeros of $\zeta(s)$ at $s = -2,-4,-6,\ldots$ lie outside the critical strip
(they have $\re(s) < 0$). Since the eta series only converges for $\re(s) > 0$, and we
restrict our analysis to $0 < \beta < 1$, the trivial zeros do not appear in our
framework. Furthermore, the prefactor $(1-2^{1-s})$ does not vanish at the trivial zeros
(since $1-s = 3,5,7,\ldots$ gives $2^{1-s} = 2^3,2^5,\ldots \ne 1$).
\end{remark}

\paragraph{Why zeros of $\mathcal{L}$ correspond to DQPTs.}
A DQPT is defined as a non-analyticity (cusp, discontinuity, or divergence) in the rate
function $\mathcal{F}(\beta,t)$, not merely a zero of the order parameter
$\mathcal{L}(\beta,t)$. We show these are equivalent in this framework.

Consider the rate function $\mathcal{F}_1(\beta,t) = -(1/\ln N)\ln|\mathcal{L}_N(\beta,t)|$
as $N \to \infty$:
\begin{enumerate}
  \item[(i)] For $(\beta,t)$ such that $\eta(\beta+it) \ne 0$:
    the numerator $\eta_N(s) \to \eta(s) \ne 0$, a bounded nonzero analytic function
    of~$s$. Write $\eta_N(s) = \eta(s) + O(N^{-\beta})$ by the Dirichlet test remainder
    bound. The denominator $Z_N(\beta) = N^{1-\beta}/(1-\beta) + O(1)$ by
    Euler--Maclaurin~\cite{titchmarsh1986}. Hence
    \begin{equation}
      |\mathcal{L}_N| = |\eta(s)|\,(1-\beta)\,N^{\beta-1}\,
      \bigl(1 + O(N^{-\min(\beta,1-\beta)})\bigr).
    \end{equation}
    The leading behavior gives $\mathcal{F}_1 \to (1-\beta)$, a smooth function.
  \item[(ii)] At a zero $s_0 = \beta_0 + it_0$ of $\eta$ (equivalently, of $\zeta$):
    write $\eta(s) = \eta'(s_0)(s-s_0) + O(|s-s_0|^2)$ near~$s_0$ (assuming simple
    zero\footnote{For zeros of multiplicity~$m$, see Remark~\ref{rem:mult}.}). At
    $s = s_0$ exactly: $\eta_N(s_0) = O(N^{-\beta_0})$ by the Dirichlet test remainder
    bound (since the full series converges to $\eta(s_0) = 0$). Hence
    $|\mathcal{L}_N(\beta_0,t_0)| = O(N^{\beta-\beta_0-1})$ and
    $\mathcal{F}_1 \to (\beta_0+1-\beta)$.
\end{enumerate}
The rate function jumps discontinuously from $(1-\beta)$ to $(\beta_0+1-\beta)$ at the
zero, a discontinuity of magnitude~$\beta_0$. Since $\beta_0 > 0$ for all nontrivial
zeros, this is a genuine non-analyticity---a first-order DQPT in the classification of
Heyl et~al.~\cite{heyl2013}.

More precisely, the derivative $d\mathcal{F}_1/dt$ diverges at $t = t_0$ (the rate
function has a cusp), because $|\mathcal{L}_N|$ passes through a minimum of order
$N^{\beta-\beta_0-1}$ and the logarithm amplifies this into a peak of height
$O(\ln N)$ in $\mathcal{F}_1$. In the thermodynamic limit, this becomes a genuine
singularity.

Therefore: zeros of $\mathcal{L}$ correspond one-to-one with DQPTs (singularities of
$\mathcal{F}_1$). This completes the heuristic argument.

\begin{remark}[Rigor]\label{rem:rigor}
The above argument is carried out at the level of formal asymptotics. A fully rigorous
statement requires establishing (i)~the existence of the thermodynamic limit
$N \to \infty$ for $\mathcal{F}_1$, with uniform convergence in $(\beta,t)$ on compact
subsets of $(0,1) \times \R$ away from zeros, and (ii)~that the singularity structure
persists in the limit. We verify both conditions numerically (Section~\ref{sec:validation}).
\end{remark}

\begin{remark}[Lehmer phenomenon]\label{rem:lehmer}
Occasionally, two zeta zeros lie extremely close together. A notable early example is
the pair near $t \approx 7005$ (gap approximately~0.04), identified by
Lehmer~\cite{lehmer1956}. Much smaller gaps ($< 10^{-3}$) are known at larger
heights~\cite{odlyzko1989}. Our zero-detection algorithm (Section~\ref{sec:zero-detection})
accounts for this with adaptive $t$-refinement.
\end{remark}

\paragraph{Catastrophic cancellation at large $t$.}
For large~$t$, the summand $n^{-(\beta+it)} = n^{-\beta}e^{-it\ln n}$ involves rapidly
oscillating phases. The partial sum $S_N(s) = \sum_{n=1}^{N} n^{-s}$ exhibits extensive
cancellation. By the variance heuristic (treating phases as approximately random), the
expected magnitude is $O\bigl((\sum_{n=1}^{N} n^{-2\beta})^{1/2}\bigr)$. For
$\beta = \tfrac{1}{2}$, this gives $O(\sqrt{\ln N})$; for $\beta > \tfrac{1}{2}$,
$O(1)$; for $\beta < \tfrac{1}{2}$, $O(N^{1/2-\beta})$. More rigorous bounds follow
from van der Corput estimates~\cite{titchmarsh1986}:
$|S_N(\beta+it)| = O(N^{1-\beta}t^{-1/2} + t^{1/2})$ for $t > 0$.

The practical consequence: with $N$ terms each contributing $O(\eps_{\mathrm{mach}})$
rounding error, the accumulated absolute error is $O(N\eps_{\mathrm{mach}})$ without
compensation (or $O(\eps_{\mathrm{mach}})$ with Kahan summation~\cite{kahan1965}). We
switch to the Riemann--Siegel formula for $t > 10^3$
(Section~\ref{sec:rs}).

\begin{remark}[Multiplicity and Robustness of Detection]\label{rem:mult}
If a zeta zero $\rho_0$ has multiplicity $m > 1$ (widely conjectured not to occur, but
not proven), two effects must be distinguished:
\begin{enumerate}
  \item[(i)] The \emph{infinite series} $\eta(s)$ vanishes to order~$m$ at~$\rho_0$:
    $\eta(s) = c_m(s-\rho_0)^m + O(|s-\rho_0|^{m+1})$. This determines the \emph{width} of the dip in
    $|\eta(s)|$ around~$\rho_0$.
  \item[(ii)] The \emph{partial sum} $\eta_N(\rho_0) = O(N^{-\beta_0})$ by the Dirichlet
    test, \emph{independently of~$m$}.
\end{enumerate}
Consequently, the rate function jump is $\beta_0$, independent of~$m$. However, zeros 
of even multiplicity do not produce a sign change in $Z_H(t)$. To ensure a 
scientifically rigorous search that is robust against such cases (or GUE violations), 
our algorithm monitors local minima of the magnitude $|\mathcal{L}(\beta, t)|$ in 
addition to sign changes of $\mathcal{G}(\beta, t)$.
\end{remark}

\subsubsection{System~2: Generalized Loschmidt Amplitude}
\label{sec:system2}

The generalized Loschmidt amplitude is defined as
\begin{equation}\label{eq:G-def}
  \mathcal{G}(\beta,t) = \frac{1}{2Z(\beta)}
  \Bigl[e^{i\vartheta(t)}\,S_N(\beta+it)
  + e^{-i\vartheta(t)}\,S_N(\beta-it)\Bigr],
\end{equation}
where $\vartheta(t)$ is the Riemann--Siegel theta function:
\begin{equation}\label{eq:theta-def}
  \vartheta(t) = \im\!\bigl(\ln\Gamma(\tfrac{1}{4}+\tfrac{it}{2})\bigr)
  - \tfrac{t}{2}\ln\pi.
\end{equation}

Since $S_N(\beta-it) = \overline{S_N(\beta+it)}$ for real~$\beta$, we have
\begin{equation}\label{eq:G-real}
  \mathcal{G}(\beta,t) = \frac{1}{Z(\beta)}\,
  \re\!\bigl[e^{i\vartheta(t)}\,S_N(\beta+it)\bigr].
\end{equation}
While $Z_H(t)$ is exactly real in the limit $N \to \infty$ due to the functional 
equation, for finite~$N$, the partial sum $S_N$ does not satisfy the 
$\zeta(s) \leftrightarrow \zeta(1-s)$ symmetry. Consequently, 
$\mathcal{G}_N(\beta, t)$ is only an approximation whose deviations from 
the real axis (or the target function) scale as the tail of the Dirichlet series, 
typically $O(N^{-\beta})$. At $\beta=1/2$, this implies a remainder of 
$O(N^{-1/2})$, which must be accounted for in the validation tolerances.


\paragraph{Divergence of $Z(\beta)$ in System~2.}
The same divergence issue arises as in System~1. The resolution is identical: we work
with the finite-$N$ quantity $\mathcal{G}_N(\beta,t)$ and extract the rate function
$\mathcal{F}_2(\beta,t) = -\lim_{N\to\infty}(1/\ln N)\ln|\mathcal{G}_N(\beta,t)|$,
which has the same behavior as $\mathcal{F}_1$: singularities at zeta zeros, smooth
elsewhere.

\paragraph{Connection to the Hardy $Z$-function.}
At $\beta = \tfrac{1}{2}$, in the limit $N \to \infty$, the partial sum
$S_N(\tfrac{1}{2}+it)$ converges to $\zeta(\tfrac{1}{2}+it)$. The Riemann--Siegel
theta function is chosen precisely so that
\begin{equation}\label{eq:hardy-z}
  Z_H(t) = e^{i\vartheta(t)}\,\zeta\!\bigl(\tfrac{1}{2}+it\bigr)
\end{equation}
is real-valued for all real~$t$~\cite{hardy1914}. This is the Hardy $Z$-function. Its
sign changes correspond one-to-one with simple zeros of $\zeta(\tfrac{1}{2}+it)$ on the
critical line.

\paragraph{Why zeros of $\mathcal{G}$ correspond to zeta zeros.}
As with System~1, $Z_N(\tfrac{1}{2})$ diverges as $N \to \infty$, so
$\mathcal{G}_N(\tfrac{1}{2},t)$ itself tends to zero. The physically meaningful
statement concerns the \emph{ratio structure}: the numerator
$\re[e^{i\vartheta(t)}S_N(\tfrac{1}{2}+it)]$ converges to $Z_H(t)$ independently of
the divergent denominator. The zeros of $\mathcal{G}_N$ (detected at finite~$N$)
correspond to the zeros of $Z_H(t)$ in the limit. Since
$Z_H(t) = 0 \iff \zeta(\tfrac{1}{2}+it) = 0$ (because $e^{i\vartheta(t)} \ne 0$),
and Hardy~\cite{hardy1914} proved infinitely many zeros lie on the critical line,
System~2 detects exactly the nontrivial zeros on the critical line.

\begin{remark}[System~2 false positives at $\beta \ne \tfrac{1}{2}$]\label{rem:false-pos}
For $\beta \ne \tfrac{1}{2}$, $\mathcal{G}(\beta,t)$ can vanish even when
$\zeta(\beta+it) \ne 0$: since $e^{i\vartheta(t)}\zeta(\beta+it)$ is generally complex
for $\beta \ne \tfrac{1}{2}$, its real part vanishes whenever the phase of
$e^{i\vartheta(t)}\zeta(\beta+it)$ equals $\pm\pi/2$. These are \emph{phase-alignment
zeros}, not zeta zeros, and constitute false positives.

\emph{Consequence for experimental design:} System~2 is reliable for zero detection
\emph{only at $\beta = \tfrac{1}{2}$}. For off-critical-line scanning
($\beta \ne \tfrac{1}{2}$), System~1 (the $\mathcal{L}$ function, based on the eta
series) must be used. The experimental design (Section~\ref{sec:protocol}) is
structured accordingly.
\end{remark}


\subsection{The RH--DQPT Equivalence}
\label{sec:equivalence}

\begin{theorem}[RH--DQPT Equivalence]\label{thm:rh-dqpt}
The Riemann Hypothesis is equivalent to the following statement: for both System~1 and
System~2, dynamical quantum phase transitions (zeros of the order parameter /
non-analyticities of the rate function) occur only at inverse temperature
$\beta = \tfrac{1}{2}$.

More precisely:
\begin{enumerate}
  \item[(a)] RH holds if and only if, for all $t \in \R$ and all $\beta \in (0,1)$ with
    $\beta \ne \tfrac{1}{2}$: $\mathcal{L}(\beta,t) \ne 0$.
  \item[(b)] At $\beta = \tfrac{1}{2}$ specifically: RH holds if and only if the zeros of
    $\mathcal{G}(\tfrac{1}{2},t)$ (as a function of~$t$) account for \emph{all}
    nontrivial zeros of~$\zeta$. (System~2 cannot be used for off-line scanning due to
    false positives; statement~(a) via System~1 is the operative equivalence.)
\end{enumerate}
\end{theorem}

\begin{proof}
\emph{Part~(a).}
($\Rightarrow$) If RH holds, all nontrivial zeros have $\re(s) = \tfrac{1}{2}$. Then
$\eta(\beta+it) = (1-2^{1-(\beta+it)})\zeta(\beta+it) \ne 0$ for $\beta \ne \tfrac{1}{2}$
in $(0,1)$ (since $\zeta$ has no zeros there and the prefactor is nonzero for
$0 < \beta < 1$ as shown in Step~3). Hence $\mathcal{L}(\beta,t) \ne 0$.

($\Leftarrow$) If RH fails, there exists $\zeta(\beta_0+it_0) = 0$ with
$\beta_0 \ne \tfrac{1}{2}$ and $0 < \beta_0 < 1$. Then $\eta(\beta_0+it_0) = 0$, so
$\mathcal{L}(\beta_0,t_0) = 0$, giving a DQPT at $\beta_0 \ne \tfrac{1}{2}$.
Contrapositive: no DQPTs off the critical line implies~RH.
\end{proof}

\paragraph{What it would mean to find a DQPT at $\beta \ne \tfrac{1}{2}$.}
This would constitute a computational disproof of the Riemann Hypothesis: the existence
of a nontrivial zero of $\zeta(s)$ with $\re(s) \ne \tfrac{1}{2}$, with immediate
consequences for analytic number theory and the distribution of primes.


\subsection{Free Energy and Singularities}
\label{sec:free-energy}

The rate function (free energy density) for System~1 is
$\mathcal{F}_1(\beta,t) = -\lim_{N\to\infty}(1/\ln N)\ln|\mathcal{L}_N(\beta,t)|$,
as derived in Section~\ref{sec:system1}, Step~5. The values are:
\begin{align}
  \mathcal{F}_1(\beta,t) &= 1-\beta && \text{(away from zeros)}, \\
  \mathcal{F}_1(\beta,t_0) &= \beta_0+1-\beta && \text{(at a zero with $\re(s_0)=\beta_0$)}, \\
  \Delta\mathcal{F}_1 &= \beta_0 && \text{(DQPT jump magnitude)}.
\end{align}
At $\beta = \beta_0 = \tfrac{1}{2}$: $\mathcal{F}_1 = 1$ at the zero,
$\mathcal{F}_1 = \tfrac{1}{2}$ away. The DQPT jump is~$\tfrac{1}{2}$.

\begin{remark}[Normalization conventions]\label{rem:norm}
Wei et~al.~\cite{wei2025} use a different normalization employing $\log_2$ rather than
$\ln$. Under their convention, $\mathcal{F}_1^{\mathrm{Wei}} = \ln 2 \cdot
\mathcal{F}_1^{\mathrm{nat}}$, giving $\mathcal{F}_1^{\mathrm{Wei}} = \ln 2$ at zeros
and $(1-\beta)\ln 2$ away. We use natural logarithms throughout; numerical values differ
by a factor of~$\ln 2$.
\end{remark}
